\section{Decision Method}
\label{sec:scoring}
\subsection{Token Probabilities}
Let $x$ be the question context and $\mathcal{C}$ a finite set of candidate
alternatives. The quoted answer text for each $c$ is represented by a token
sequence $y_c=(y_{c,1},\ldots,y_{c,m_c})$, following the model's assistant
prefix, with $m_c>0$. Tokenization preserves the boundary between context and
answer. The length $m_c$ counts the actual tokens in the scored candidate
suffix, including its quotation marks and escaped content. It excludes prompt
tokens, any reasoning trace, image tokens, and unscored end-of-turn tokens.
Let $\mathcal{V}$ be the model vocabulary
and $z_{c,t,v}$ the logit for token $v$ given $(x,y_{c,<t})$.
The probability of each supplied candidate token is
\begin{equation}
  p_{c,t}=p_\theta(y_{c,t}\mid x,y_{c,<t})
  =\frac{\exp(z_{c,t,y_{c,t}})}
         {\sum_{v\in\mathcal{V}}\exp(z_{c,t,v})}.
  \label{eq:token-probability}
\end{equation}
The denominator includes the full vocabulary. Tokens are evaluated under
teacher forcing, without sampling, temperature rescaling, or top-$k$/top-$p$
filtering. The answer text, rather than an arbitrary option label, determines
the score.

For numerical stability, the log-probability $\ell_{c,t}=\log p_{c,t}$ is
written as
\begin{equation}
  \ell_{c,t}=z_{c,t,y_{c,t}}-a_{c,t}
  -\log\sum_{v\in\mathcal{V}}\exp(z_{c,t,v}-a_{c,t}),
  \qquad a_{c,t}=\max_{v\in\mathcal{V}}z_{c,t,v}.
  \label{eq:token-log-probability}
\end{equation}

\subsection{Length-Normalized Candidate Scores}
The autoregressive chain rule gives the sequence likelihood and its logarithm:
\begin{equation}
  L_c=p_\theta(y_c\mid x)=\prod_{t=1}^{m_c}p_{c,t},
  \qquad
  r(c\mid x)=\log L_c=\sum_{t=1}^{m_c}\ell_{c,t}.
  \label{eq:sequence-likelihood}
\end{equation}
The sum accumulates a negative contribution for each token, which penalizes
longer candidates even when their mean token log-probabilities are identical.
JET therefore uses the mean token log-probability as its decision score:
\begin{equation}
  s(c\mid x)=\frac{r(c\mid x)}{m_c}
  =\frac{1}{m_c}\sum_{t=1}^{m_c}\ell_{c,t}.
  \label{eq:score}
\end{equation}
Its exponential is the geometric mean of the candidate's conditional token
probabilities, $\exp(s(c\mid x))=L_c^{1/m_c}$. A softmax over these scores,
with no additional temperature scaling ($T=1$), gives normalized decision
weights:
\begin{equation}
  q(c\mid x,\mathcal{C})
  =\frac{L_c^{1/m_c}}{\sum_{d\in\mathcal{C}}L_d^{1/m_d}}
  =\frac{\exp(s(c\mid x)-M)}
        {\sum_{d\in\mathcal{C}}\exp(s(d\mid x)-M)},
  \quad M=\max_{d\in\mathcal{C}}s(d\mid x).
  \label{eq:probability}
\end{equation}
Equations~\ref{eq:token-probability} and~\ref{eq:probability} normalize over
different sets: vocabulary tokens and candidate scores, respectively. Two
candidates with lengths 2 and 10 and mean token log-probability $-0.2$ both
receive weight $0.5$, despite raw log-likelihoods of $-0.4$ and $-2.0$.
The weights sum to one, but are not the candidate-restricted sequence
likelihoods $L_c/\sum_d L_d$ or calibrated probabilities of correctness.
The API retains its probability fields for these weights.

Length normalization reduces the direct cumulative length penalty; it does not
make scores invariant to wording, tokenization, or answer priors. Appending
easily predicted text can raise the mean score, and scores are not summed over
alternative phrasings. Adding or duplicating a candidate changes the normalized
weights. Calibration therefore requires separate empirical evaluation.
All text and multimodal execution paths use the same mean-token scoring rule.

\subsection{Decision Rules}
For a categorical decision, JET selects
\begin{equation}
  \hat c=\operatorname*{arg\,max}_{c\in\mathcal{C}}q(c\mid x,\mathcal{C})
  =\operatorname*{arg\,max}_{c\in\mathcal{C}}s(c\mid x),
  \label{eq:decision}
\end{equation}
with ties resolved by a fixed ordering. A binary question uses
$\mathcal{C}=\{\mathrm{false},\mathrm{true}\}$ and returns
$q(\mathrm{true}\mid x,\mathcal{C})$, the true candidate's decision weight.
An ordinal question with $K$ ordered levels returns the weighted index,
$\sum_{i=0}^{K-1} i\,q(i\mid x,\mathcal{C})$, which may be fractional.
These rules require no additional classification head or task-specific training.

Candidate labels are distinct from answer text. Different labels with identical
text receive the same score and use the fixed tie-breaking rule; binary
and ordinal answer descriptions must be distinct. In evaluation, any label with
the gold answer text is accepted, and its weight is summed with those of
equivalent labels when computing gold negative log-likelihood (NLL). This
diagnostic is the negative log of the gold decision weight, not the raw
sequence negative log-likelihood from Equation~\ref{eq:sequence-likelihood}.

\subsection{Optional Reasoning and Determinism}
JET can generate a reasoning trace once per question and include it in the
context shared by all candidates. Candidate evaluation remains fixed conditional
on that trace, but the complete procedure then includes sampling. The main
execution comparisons disable reasoning generation.

Without reasoning generation, the decision rule is deterministic for fixed
model outputs. Floating-point differences across devices and execution settings
can still change decision weights and decisions. Section~\ref{sec:repeatability}
describes the numerical validation boundary; agreement in a tested configuration
does not imply bitwise reproducibility across hardware.
